\section{Experimental Setup}
\label{sec:setup}

\subsection{States and District Counts}

We evaluate all eight algorithms on three states that span a representative
range of congressional district counts and geographic character:

\begin{itemize}
  \item \textbf{North Carolina (NC)}, $k = 14$.
        A contested state with a history of contested congressional maps.
        Moderately large tract graph ($N \approx 2200$ census tracts),
        county structure that frequently produces split counties, and a
        competitive partisan baseline.
        $k = 14$ is a regime in which all single-scale chain methods mix
        adequately; NC serves primarily as the baseline state.

  \item \textbf{Wisconsin (WI)}, $k = 8$.
        A small-$k$ state where mixing is straightforward for all methods.
        WI highlights runtime differences and the overhead costs of
        algorithms designed for larger problems (multi-scale, VRA boosts).
        It also provides a clean test of \flip{}: with $k = 8$, individual
        boundary tracts constitute a larger fraction of each district's
        perimeter than for larger $k$.

  \item \textbf{Texas (TX)}, $k = 38$.
        The canonical large-$k$ test case for redistricting ensembles.
        TX has $N \approx 5200$ census tracts, a complex county structure
        ($\approx 100$ active county regions in the tract graph), and a
        plan space that is too large for single-scale chains to explore
        adequately without multi-scale acceleration.
        TX is the primary state for autocorrelation comparisons.
\end{itemize}

\subsection{Common Configuration}

All runs use:
\begin{itemize}
  \item \textbf{Base seed}: 42 (set via \texttt{--seed 42}).
  \item \textbf{Population tolerance}: $\epsilon = 0.5\%$ (the default
        balance tolerance for the compositor).
  \item \textbf{Structure}: \texttt{ratio-optimal} (GeoSection bisection)
        for all methods, providing the same initial plan as the starting
        point.
  \item \textbf{Weights}: \texttt{geographic} (default edge weights,
        no county or VRA signal in the underlying graph).
  \item \textbf{Hardware}: Intel Core i7-11800H (8 cores, 16 threads,
        2.3~GHz base), 32~GB DDR4-3200. All runs use a single thread to
        ensure runtime comparability.
  \item \textbf{Census year}: 2020.
\end{itemize}

\paragraph{Chain-method step budget.}
For chain methods (\sburst{}, \flip{}, \fr{}, \msp{}, \ms{},
\sburstf{}, \vra{}), we run $T = 1{,}000$ steps per state.
This budget is deliberately modest; the goal is not to produce a
production-quality ensemble but to compare algorithms under the same
compute constraint.
For TX, 1000 steps means the chain touches each district approximately
$1000 \times (2/38) \approx 53$ times in expectation---adequate for
comparing autocorrelation structure but insufficient for a litigation-quality
ensemble (which would require $T \geq 50{,}000$).

\paragraph{SMC particle budget.}
For \smc{}, we use $N = 1{,}000$ particles, matching the step budget of
the chain methods in terms of total plan evaluations.
The SMC run is a single pass (no MCMC warm-up chains); particles are
resampled based on population-balance weights.

\subsection{Metrics}

We report four metrics for each algorithm--state pair.

\paragraph{Edge-cut score (EC).}
Let $G = (V, E)$ be the tract adjacency graph and $\pi : V \to [k]$
a district assignment.
The edge-cut score is:
\[
  \ec(\pi) = |\{(u,v) \in E : \pi(u) \neq \pi(v)\}|.
\]
Lower \ec{} means fewer cross-district adjacencies---a proxy for
compactness.
We report the \ec{} of the final plan produced by each algorithm
(for chain methods, the plan at step $T$).

\paragraph{Mean Polsby-Popper (PP).}
For district $d$ with area $A_d$ and perimeter $P_d$, the Polsby-Popper
score is:
\[
  \mathrm{PP}_d = \frac{4\pi A_d}{P_d^2}.
\]
We report the mean $\pp = \frac{1}{k} \sum_{d=1}^k \mathrm{PP}_d$ of the
final plan.
Higher \pp{} means more compact districts (closer to a circle).
\pp{} is computed from the tract-level geometries projected to district
boundaries.

\paragraph{Partisan seat stability.}
We run each algorithm five times with seeds $\{42, 43, 44, 45, 46\}$ and
report the fraction of runs that produce the same Democratic seat count
under the 2020 presidential vote totals.
High stability indicates that the algorithm reliably produces plans with a
particular partisan outcome; low stability indicates sensitivity to
seed.
Note: partisan seat stability measures algorithmic consistency, not
algorithmic fairness.

\paragraph{Wall-clock runtime.}
Total wall-clock time in seconds from job start to final plan output,
including all preprocessing (graph loading, coarsening derivation, initial
plan construction) but excluding census data download.
Measured on the single-threaded i7-11800H using the
\texttt{--workers~1} flag.

\paragraph{Lag-100 autocorrelation (chain methods only).}
For chain-based methods, we record the \ec{} statistic at each step and
compute the sample autocorrelation at lag $\ell = 100$:
\[
  \acf(100) = \frac{\sum_{t=1}^{T-100} (\ec_t - \bar{\ec})
                    (\ec_{t+100} - \bar{\ec})}
                   {\sum_{t=1}^{T} (\ec_t - \bar{\ec})^2}.
\]
Lower $\acf(100)$ indicates faster mixing: samples 100 steps apart are
more nearly independent.

\subsection{Single-Run Caveat}
\label{sec:caveat}

\textbf{All compactness and runtime figures in this paper are single-run
results (seed 42) and should be interpreted as illustrative, not definitive.}
Ensemble estimates with meaningful confidence intervals require $R \geq 10$
independent runs, which is outside the scope of this comparative paper.
The figures presented here are consistent with repeated informal testing
but have not been subjected to formal statistical inference.
The decision tree in Section~\ref{sec:decision} is based on the
structural properties of each algorithm (as established in G.6--G.13),
not solely on the single-run empirical results of this paper.

\subsection{Algorithm--CLI Mapping}

Table~\ref{tab:cli} gives the \texttt{--search} flag that invokes each
algorithm in the \texttt{redist} CLI.

\begin{table}[ht]
\centering
\caption{CLI flags for each algorithm.}
\label{tab:cli}
\begin{tabular}{@{}lll@{}}
\toprule
Algorithm & \texttt{--search} flag & Paper \\
\midrule
Short-Burst     & \texttt{multi}                 & G.6 \\
SMC             & (standalone mode)              & G.7 \\
Flip            & \texttt{flip}                  & G.8 \\
Forest ReCom    & \texttt{forest-recom}          & G.9 \\
Merge-Split     & \texttt{merge-split}           & G.10 \\
Multi-scale     & \texttt{multiscale}            & G.11 \\
ShortBurstForest & \texttt{shortburst-forest}   & G.12 \\
VRA-Aware ReCom & \texttt{vra-recom}             & G.13 \\
\bottomrule
\end{tabular}
\end{table}
